\documentclass[11pt,a4paper]{article}
\usepackage[utf8]{inputenc}
\usepackage[T1]{fontenc}
\usepackage[ngerman]{babel}
\usepackage{geometry}
\usepackage{fancyhdr}
\usepackage{amsmath}
\usepackage{amssymb}

\geometry{margin=25mm}
\pagestyle{fancy}
\fancyhf{}
\renewcommand{\headrulewidth}{0.4pt}
\lhead{\textbf{Name:} \studentname}
\rhead{\textbf{Matr.-Nr.:} \studentnumber}
\cfoot{\thepage}

\newcommand{\studentname}{Matthias Jan\ss{}en}
\newcommand{\studentnumber}{1871808}
\newcommand{\course}{Grundlagen der Computergrafik}
\newcommand{\institute}{Universit\"at Trier}
\newcommand{\sheetnumber}{06}
\newcommand{\tutor}{Jan Niclas Ruppenthal\\Prof. Dr. Stephan Diehl}

\title{\course}
\author{\studentname\\Matrikelnummer: \studentnumber\\\institute}
\date{\today}

\begin{document}

\maketitle

\begin{center}
    \textbf{Abgabe:} \"Ubungsblatt \sheetnumber \\
    \vspace{4mm}
    \textbf{Tutor:} \tutor
\end{center}

\bigskip

\section*{Aufgaben}

% ----------------------------------------------------------------
\section{Aufgabe 1 -- Triangle Strips/Fans}

Der K\"orper ist eine pentagonale Pyramide mit Spitze $F$ und
Grundfl\"ache $ABCDE$, bestehend aus den 5 Seitenfacetten:
$(E,A,F)$, $(A,B,F)$, $(B,C,F)$, $(C,D,F)$, $(D,E,F)$.

\subsection*{a) Triangle Strips}

Je zwei benachbarte Seitendreiecke teilen ausschlie\ss{}lich Kanten,
die $F$ enthalten. Daher kann ein Strip maximal 2 aufeinanderfolgende
Dreiecke abdecken, z.\,B.\ $(E,A,F,B)$ liefert $(E,A,F)$ und $(A,F,B)$.
F\"ur 5 Dreiecke werden mindestens \textbf{3 Triangle Strips} ben\"otigt:

\medskip
\noindent\textbf{Strip 1} (2 Dreiecke):
\[
  (E,\; A,\; F,\; B)
\]
Dreiecke: $(E,A,F)$,\; $(A,F,B)$

\noindent\textbf{Strip 2} (2 Dreiecke):
\[
  (B,\; C,\; F,\; D)
\]
Dreiecke: $(B,C,F)$,\; $(C,F,D)$

\noindent\textbf{Strip 3} (1 Dreieck):
\[
  (D,\; E,\; F)
\]
Dreiecke: $(D,E,F)$

\subsection*{b) Triangle Fans}

Minimale Anzahl: \textbf{1 Triangle Fan}

\medskip
\noindent\textbf{Fan} mit Zentrum $F$ (5 Dreiecke):
\[
  (F,\; E,\; A,\; B,\; C,\; D,\; E)
\]
Dreiecke: $(F,E,A)$,\; $(F,A,B)$,\; $(F,B,C)$,\; $(F,C,D)$,\; $(F,D,E)$

% ----------------------------------------------------------------
\section{Aufgabe 2 -- Polygon}

\subsection*{a) Innenwinkelsumme eines $n$-seitigen konvexen Polygons}

\textbf{Beh.:} $\displaystyle\sum_{i=1}^{n} \beta_i = (n-2)\cdot\pi$

\medskip
\noindent\textbf{I.A.} ($n=3$): $\quad\pi = (3-2)\cdot\pi$ \checkmark

\medskip
\noindent\textbf{I.V.}: F\"ur ein $n$-seitiges Polygon gilt $\sum \beta_i = (n-2)\cdot\pi$.

\medskip
\noindent\textbf{I.S.} ($n \to n+1$): Eine Diagonale zerlegt das $(n{+}1)$-seitige
Polygon in zwei Polygone mit $n_L$ und $n_R$ Ecken, wobei $n_L + n_R = n + 3$.
Mit der I.V. folgt:
\[
  \sum \beta_i
  = (n_L - 2)\pi + (n_R - 2)\pi
  = (n_L + n_R - 4)\pi
  = (n - 1)\pi
  = \bigl((n+1)-2\bigr)\pi. \qquad\square
\]

\subsection*{b) Summe der Au\ss{}enwinkel}

Da $\alpha_i = \pi - \beta_i$ an jedem Eckpunkt:
\[
  \sum_{i=1}^{n} \alpha_i
  = n\pi - (n-2)\pi
  = 2\pi. \qquad\square
\]

\end{document}